\documentclass{article}
\usepackage{amsmath, amssymb, amsthm}
\usepackage{hyperref}
\usepackage{geometry}
\geometry{margin=1in}

\title{Erd\H{o}s Problem \#195}
\date{Last updated: 2025-08-31}
\author{Source: erdosproblems.com}

\begin{document}
\maketitle

\noindent\textbf{Status:} OPEN \quad \textbf{No prize}

\noindent\textbf{Tags:} arithmetic progressions

\medskip

\noindent\textbf{Problem Statement:}

What is the largest $k$ such that in any permutation of $\mathbb{Z}$ there must exist a monotone $k$-term arithmetic progression $x_1&#60;\cdots&#60;x_k$?




Geneson \cite{Ge19} proved that $k\leq 5$. Adenwalla \cite{Ad22} proved that $k\leq 4$. See also [194] and [196] .




References

[Ad22] Adenwalla, S., Avoiding Monotone Arithmetic Progressions in Permutations of Integers . arXiv:2211.04451 (2022).

[Ge19] Geneson, Jesse, Forbidden arithmetic progressions in permutations of subsets of the integers . Discrete Math. (2019), 1489-1491.


\bigskip
\noindent\small{Source: \url{https://www.erdosproblems.com/195}}
\end{document}
